\section{Conclusion}
\label{sec:conclusion}

Under one corpus, one evaluator, and one measurement protocol, no universal
winner emerged among the twelve finite-budget states (RQ1): RNDVoC
led PESQ and STOI, HiFi-GAN V1 led UTMOS, Vocos led RTF on both lanes,
HiFi-GAN V2 led footprint, and RFWave held the intermediate point.  Deployment transformations were conditional on architecture, operator
graph, and backend (RQ2): FP16 casting was most consistently beneficial,
quantization and compilation were heterogeneous, static ONNX INT8 coupled
large size reductions to material quality loss, dense masks measured
tolerance to zeroed weights rather than sparse execution, and RFWave
solver reduction gave the one monotone speed-quality control.  At the moderate quality, parameter, and storage
vector the artifact-level selector (RQ3) chose FP16-cast Vocos on B200
and dynamic-INT8 Vocos on CPU; the tighter vector selected untransformed
Vocos on CPU, and RNDVoC stayed the highest-PESQ feasible state at both
vectors, never the RTF-minimizing choice.  These conclusions are bounded
by proxy metrics, an adaptive single-speaker evaluation set, single
trajectories, and non-host-paired timing.

\paragraph{Availability.}
The source code of both software packages, the complete experimental
log, and the admission, exclusion, and result registers under
\texttt{artifacts/} are public at
\url{https://github.com/rahul1-bot/VOCODE-Efficient-Neural-Vocoders}; the
graded state is the latest commit on or before 31 July 2026.  LJSpeech is
public under its released license, and resolved run configurations record
the environment.
